% ---------------------------------------------------------------------------
%  Preamble of the book. The typography follows Evan Chen's "An Infinitely
%  Large Napkin": Latin Modern 11pt, bold sans-serif headings, maroon links and
%  contents, crimson chapter numerals and section signs, a blue rounded box for
%  the most important items, a rust box for examples, a green bar for remarks.
%  Book-specific settings live in bookmeta.tex; extra boxes and tools in kit.tex.
% ---------------------------------------------------------------------------
\documentclass[paper=letter,fontsize=11pt,twoside,open=any,headings=normal,
               numbers=noenddot]{scrbook}

\usepackage[T1]{fontenc}
\usepackage[utf8]{inputenc}
\usepackage{lmodern}
\usepackage[letterpaper,left=1.3in,right=1.3in,top=1.15in,bottom=1.1in]{geometry}
\usepackage{amsmath,amssymb}
\usepackage{microtype}
\usepackage{xurl}
\usepackage{mathtools,bm}
\usepackage{booktabs,array,tabularx,longtable,multirow,makecell,colortbl}
\usepackage{enumitem}
\usepackage{pifont}
\usepackage{etoolbox}
\usepackage{needspace}
\usepackage{xspace}
\usepackage{xcolor}
\usepackage{graphicx}
\usepackage[font=small,labelfont={sf,bf},labelsep=period,justification=raggedright,singlelinecheck=false]{caption}
\usepackage{tikz}
\usetikzlibrary{arrows.meta,positioning,shapes.geometric,shapes.misc,calc,fit,backgrounds,decorations.pathreplacing,decorations.pathmorphing,matrix}
\usepackage{pgfplots}
\pgfplotsset{compat=1.17}
\usepackage{siunitx}
\usepackage{listings}
\usepackage[most]{tcolorbox}
\usepackage[backend=biber,style=alphabetic,maxalphanames=3,minalphanames=3,sorting=anyt,
            maxbibnames=6,minbibnames=3,giveninits=true,doi=false,isbn=false,url=true,eprint=true,
            backref=true,urldate=long]{biblatex}
\DefineBibliographyStrings{english}{backrefpage={p.\addnbspace},backrefpages={pp.\addnbspace}}
\DeclareFieldFormat{labelalpha}{\mdseries #1}  % citation labels stay in medium weight inside bold table cells
\renewcommand*{\finentrypunct}{}
\renewcommand*{\mkbibnamefamily}[1]{\textsc{#1}}  % family names in small caps, as in the Napkin's bibliography

\addbibresource{refs_auto.bib}
\addbibresource{refs_manual.bib}

% \fillin{...}: a visible red marker for text that still has to be written. tools/check.sh fails while any remain.
\newcommand{\fillin}[1]{\textcolor{red}{\textbf{[FILL IN: #1]}}}
% ---------------------------------------------------------------------------
%  Book-level settings. Edit these first; everything else reads them.
%  scripts/new_book.sh fills in the title, short title and date; \fillin{...}
%  prints a red marker that tools/check.sh reports until it is replaced.
% ---------------------------------------------------------------------------
\newcommand{\booktitle}{BOOK TITLE}
\newcommand{\booksubtitle}{\fillin{subtitle, e.g.\ a reading guide, with a crash course in X}}
\newcommand{\bookshorttitle}{SHORT TITLE}            % running head on even pages, and the cover sticker
\newcommand{\bookcoverline}{\fillin{cover line}}      % small line under the title on the cover, e.g. a reading guide
\newcommand{\bookauthor}{Compiled with Claude}
\newcommand{\bookdate}{BOOK DATE}                     % American style: November 3, 2026
\newcommand{\bookversion}{v1.0.BOOKSTAMP}             % bump on every delivered rebuild
\newcommand{\booksubject}{BOOK TITLE}                 % PDF metadata
\newcommand{\bookkeywords}{}                          % PDF metadata, comma-separated
% Title of the "connection" box (kit.tex): it links a section to the book's main application or to a neighboring
% field, e.g. "Robotics connection" or "Clinical connection". If every chapter is already about the main topic,
% name it for what it points to instead, e.g. "Related idea" or "In practice".
\newcommand{\connectionname}{Connection}


\usepackage[colorlinks=true,linkcolor=maroon,citecolor=citegreen,urlcolor=blue,
            pdftitle={\booktitle},pdfauthor={\bookauthor},pdfsubject={\booksubject},pdfkeywords={\bookkeywords},
            bookmarksnumbered=true,breaklinks=true]{hyperref}
\usepackage{bookmark}
\usepackage[nameinlink,noabbrev]{cleveref}

% ---- colors (the Napkin's palette) ------------------------------------------
\definecolor{maroon}{HTML}{800000}      % links, contents, headings accents
\definecolor{crimson}{HTML}{BF0040}     % chapter numerals, section signs
\definecolor{citegreen}{RGB}{0,128,0}   % citations
\definecolor{propborder}{HTML}{0000FF}  % "Proposition" box
\definecolor{propfill}{HTML}{F4FBFA}
\definecolor{proptitle}{HTML}{07608F}
\definecolor{exborder}{HTML}{94340E}    % "Example" box
\definecolor{exfill}{HTML}{FEF8F7}
\definecolor{rembar}{HTML}{159B52}      % "Remark" box
\definecolor{remtitle}{HTML}{3D8252}
\definecolor{remfill}{HTML}{F2F9F2}
% names used by the figures
\colorlet{navy}{proptitle}
\colorlet{teal}{remtitle}
\colorlet{amber}{exborder}
\colorlet{rose}{crimson}
\definecolor{slate}{HTML}{595959}
\definecolor{mist}{HTML}{F4FBFA}
\definecolor{sand}{HTML}{FEF3EC}
\definecolor{line}{HTML}{C8C8C8}
% reading tiers: A essential (blue box), B important (rust box), C context (green bar)
\colorlet{tierA}{proptitle}
\colorlet{tierB}{exborder}
\colorlet{tierC}{remtitle}

% ---- text layout -----------------------------------------------------------
\setcounter{tocdepth}{2}
\setcounter{secnumdepth}{3}
\setlength{\parindent}{11.5pt}
\setlength{\parskip}{0pt plus 0.5pt}
\raggedbottom
\makeatletter
\setlength{\@fptop}{0pt}\setlength{\@fpsep}{14pt plus 1fil}\setlength{\@fpbot}{0pt plus 1fil}
\makeatother
\renewcommand{\topfraction}{0.85}\renewcommand{\bottomfraction}{0.75}
\renewcommand{\textfraction}{0.12}\renewcommand{\floatpagefraction}{0.85}
\setlength{\emergencystretch}{2em}
\clubpenalty=10000
\widowpenalty=10000
\setlist{nosep,leftmargin=1.8em,itemsep=3pt,topsep=4pt}
\renewcommand{\arraystretch}{1.1}
\let\cleardoublepage\clearpage  % no blank verso pages (\mainmatter and \frontmatter use \cleardoubleoddpage)
\let\cleardoubleoddpage\clearpage

% ---- headings (KOMA), after the Napkin ---------------------------------------
\renewcommand*{\thesubsection}{\thesection.\roman{subsection}}
\renewcommand*{\thesubsubsection}{\thesubsection.\alph{subsubsection}}
\setkomafont{disposition}{\sffamily\bfseries}
\setkomafont{part}{\normalfont\bfseries\fontsize{24}{28}\selectfont\color{black}}
\setkomafont{partnumber}{\normalfont\bfseries\fontsize{56}{60}\selectfont\color{crimson}}
\setkomafont{chapter}{\sffamily\bfseries\fontsize{24}{28}\selectfont}
\setkomafont{section}{\sffamily\bfseries\Large}
\setkomafont{subsection}{\sffamily\bfseries\normalsize}
\setkomafont{subsubsection}{\sffamily\bfseries\normalsize}
\setkomafont{paragraph}{\normalfont\bfseries}
\setkomafont{pageheadfoot}{\normalfont\small}
\setkomafont{pagenumber}{\normalfont\small}
\renewcommand*{\partformat}{\thepart}
\renewcommand*{\raggedpart}{\centering}
\renewcommand*{\partheadstartvskip}{\vspace*{0.24\textheight}}
\renewcommand*{\partheadmidvskip}{\par\nobreak\vspace{0.9\baselineskip}}
\renewcommand*{\chapterformat}{{\color{crimson}\fontsize{54}{54}\selectfont\thechapter}\hspace{0.45em}}
\renewcommand*{\chapterheadstartvskip}{\vspace*{-0.5\baselineskip}}
\renewcommand*{\chapterheadendvskip}{\vspace{1.5\baselineskip}}
\renewcommand*{\sectionformat}{\textcolor{crimson}{\S\thesection}\enskip}
\renewcommand*{\subsectionformat}{\textcolor{crimson}{\S\thesubsection}\enskip}
\renewcommand*{\subsubsectionformat}{\textcolor{crimson}{\S\thesubsubsection}\enskip}
\RedeclareSectionCommand[beforeskip=2.3\baselineskip plus 4pt minus 2pt,afterskip=0.75\baselineskip plus 2pt]{section}
\RedeclareSectionCommand[beforeskip=1.6\baselineskip plus 3pt minus 2pt,afterskip=0.55\baselineskip plus 1pt]{subsection}
\RedeclareSectionCommand[beforeskip=1.2\baselineskip plus 2pt minus 1pt,afterskip=0.4\baselineskip]{subsubsection}
\RedeclareSectionCommand[indent=0pt,beforeskip=0.9\baselineskip plus 2pt minus 1pt,
                         afterskip=-0.5em]{paragraph}

% ---- running heads and feet --------------------------------------------------
\usepackage[headwidth=\textwidth,footwidth=\textwidth]{scrlayer-scrpage}
\pagestyle{scrheadings}
\clearpairofpagestyles
\automark[chapter]{chapter}
\renewcommand*{\chaptermarkformat}{\thechapter\enskip}
\lehead{\pagemark}
\rehead{\textsf{\textbf{\bookshorttitle}}, compiled with \emph{Claude} (\bookversion)}
\lohead{\headmark}
\rohead{\pagemark}
\cfoot[\pagemark]{}
\KOMAoptions{headsepline=0.6pt}

% keep section headings with the box that follows them
\pretocmd{\section}{\Needspace{9\baselineskip}}{}{\PackageWarning{preamble}{could not patch \string\section}}
\pretocmd{\subsection}{\Needspace{7\baselineskip}}{}{\PackageWarning{preamble}{could not patch \string\subsection}}

% ---- table of contents: maroon entries, black page numbers ------------------
\makeatletter\renewcommand*{\@pnumwidth}{2.3em}\makeatother  % room for three bold digits
\DeclareTOCStyleEntry[indent=0pt,numwidth=46pt,beforeskip=1.5em plus 2pt,linefill=\hfill,
                      entryformat=\sffamily\bfseries\large\color{maroon},
                      pagenumberformat=\sffamily\bfseries\large\color{black}]{tocline}{part}
\DeclareTOCStyleEntry[indent=0pt,numwidth=46pt,beforeskip=0.75em plus 1pt,linefill=\hfill,
                      entryformat=\sffamily\bfseries\color{maroon},
                      pagenumberformat=\sffamily\bfseries\color{black}]{tocline}{chapter}
\DeclareTOCStyleEntry[indent=17pt,numwidth=35pt,beforeskip=0pt,
                      entryformat=\color{maroon},
                      pagenumberformat=\color{black}]{tocline}{section}

\providecommand{\npapers}{0}
% AUTO-GENERATED by tools/make_papers_listing.py -- do not edit
\renewcommand{\npapers}{0}
   % generated by tools/make_papers_listing.py; must exist so that latexmk tracks it

% ---- tiers / paper database -----------------------------------------------
% AUTO-GENERATED by tools/build_meta.py -- do not edit

                         % generated by tools/build_meta.py (do not edit)

\csdef{P@stars@A}{\ding{72}\ding{72}\ding{72}}
\csdef{P@stars@B}{\ding{72}\ding{72}}
\csdef{P@stars@C}{\ding{72}}
\csdef{P@word@A}{ESSENTIAL}
\csdef{P@word@B}{IMPORTANT}
\csdef{P@word@C}{CONTEXT}

\newcommand{\pfield}[2]{\csname P@#1@#2\endcsname}              % \pfield{title}{key}
\newcommand{\ptier}[1]{\csname P@tier@#1\endcsname}
\newcommand{\pcolor}[1]{pc@#1}
\newcommand{\pstars}[1]{\csname P@stars@\csname P@tier@#1\endcsname\endcsname}
\newcommand{\pword}[1]{\csname P@word@\csname P@tier@#1\endcsname\endcsname}
\newcommand{\arxivlink}[1]{\href{https://arxiv.org/abs/\pfield{arx}{#1}}{arXiv:\pfield{arx}{#1}}}

% Card label (run-in heading inside a paper card), in the color of the box
\colorlet{cardcolor}{tierA}
\newcommand{\cl}[1]{\par\noindent\textsf{\textbf{\color{cardcolor}#1}}\enspace}

\newcommand{\metaline}[1]{%
  \pfield{auth}{#1}%
  \ \textperiodcentered\ \pfield{date}{#1}%
  \ifcsempty{P@venue@#1}{}{\ (\pfield{venue}{#1})}%
  \ifcsempty{P@org@#1}{}{\expandafter\ifx\csname P@org@#1\endcsname\csname P@auth@#1\endcsname\else\ \textperiodcentered\ \pfield{org}{#1}\fi}%
  \ \textperiodcentered\ \arxivlink{#1}%
  \ \textperiodcentered\ \cite{#1}}

% ---- boxes -------------------------------------------------------------------
% One counter for every box, reset in each section, as in the Napkin ("Paper 5.2.1", "Example 5.2.2").
\newcounter{nbox}[section]
\renewcommand{\thenbox}{\thesection.\arabic{nbox}}

\tcbset{
  napkin/.style={enhanced jigsaw,breakable,left=10pt,right=10pt,top=6pt,bottom=7pt,
      before skip=8pt plus 2pt,after skip=10pt plus 2pt,
      before upper={\parindent=0pt\parskip=3pt\relax}},
  boxA/.style={napkin,arc=7pt,boxrule=1.1pt,colframe=propborder,colback=propfill},
  boxB/.style={napkin,sharp corners,boxrule=0.7pt,colframe=exborder,colback=exfill},
  boxC/.style={napkin,sharp corners,boxrule=0pt,frame hidden,colback=remfill,
      borderline west={2.6pt}{0pt}{rembar}},
  boxR/.style={napkin,sharp corners,boxrule=0pt,frame hidden,colback=crimson!6,
      borderline west={2.6pt}{0pt}{crimson}},
}

% little samples of the three card boxes, for the legend of the reading tiers
\newcommand{\swatchA}{\tikz[baseline=-0.55ex]{\draw[propborder,fill=propfill,line width=0.8pt,rounded corners=2pt] (0,-0.9ex) rectangle (2.6em,1.5ex);}}
\newcommand{\swatchB}{\tikz[baseline=-0.55ex]{\draw[exborder,fill=exfill,line width=0.6pt] (0,-0.9ex) rectangle (2.6em,1.5ex);}}
\newcommand{\swatchC}{\tikz[baseline=-0.55ex]{\fill[remfill] (0,-0.9ex) rectangle (2.6em,1.5ex);\fill[rembar] (0,-0.9ex) rectangle (0.28em,1.5ex);}}

% Paper card, common part: \cardbegin{tier}{bibkey}{name}{full title}{meta line}
\newcommand{\cardbegin}[5]{%
  \Needspace{7\baselineskip}%
  \begin{tcolorbox}[box#1]%
  \colorlet{cardcolor}{tier#1}%
  \refstepcounter{nbox}\label{card:#2}%
  {\color{cardcolor}\sffamily\bfseries Paper~\thenbox}{\color{cardcolor}\sffamily\ (#3)}%
  \hfill{\color{crimson}\small\csname P@stars@#1\endcsname}\par
  \vspace{1pt}%
  {\small\itshape #4\par}%
  {\footnotesize\sffamily\color{slate}#5\par}%
  \vspace{1pt}{\color{cardcolor!45}\hrule height 0.4pt}\vspace{3pt}%
}

% \begin{papercard}{bibkey}{Display name} ... \end{papercard}    (arXiv papers registered with tools/add_arxiv.py)
\NewDocumentEnvironment{papercard}{m m}{%
  \ifcsname P@title@#1\endcsname\else\PackageError{papercard}{Unknown paper key #1}{Add it with tools/add_arxiv.py}\fi
  \edef\cardtier{\csname P@tier@#1\endcsname}%
  \expandafter\cardbegin\expandafter{\cardtier}{#1}{#2}{\pfield{title}{#1}}{\metaline{#1}}%
}{\end{tcolorbox}}

% Card for sources without arXiv metadata (books, journal-only papers, blog posts, model cards):
%   \begin{webcard}{bibkey}{Display name}{Full title}{meta line}{A|B|C}     (bibkey lives in refs_manual.bib)
\NewDocumentEnvironment{webcard}{m m m m m}{%
  \cardbegin{#5}{#1}{#2}{#3}{#4\ \textperiodcentered\ \cite{#1}}%
}{\end{tcolorbox}}

% Callouts: insightbox (blue box), readingbox and cautionbox (bar boxes with a run-in title)
%   \begin{insightbox}[optional title] ... \end{insightbox}
% The \Needspace below keeps a box title from being stranded at the bottom of a page. Do NOT lower it globally to
% fix one bad page break: every box in the book moves and the whole book repaginates. Fix single pages locally
% (\enlargethispage, a shorter sentence, moving a float).
\NewDocumentCommand{\calloutstart}{m m m m +m}{% style, accent color, type, title, what follows the title
  \Needspace{6\baselineskip}%
  \begin{tcolorbox}[#1]%
  \refstepcounter{nbox}%
  {\color{#2}\sffamily\bfseries #3~\thenbox}%
  \IfBlankTF{#4}{}{{\color{#2}\sffamily\ (#4)}}%
  #5%
}
\NewDocumentEnvironment{insightbox}{O{}}{%
  \calloutstart{boxA,unbreakable}{proptitle}{Key idea}{#1}{\par\vspace{1pt}}%
}{\end{tcolorbox}}
\NewDocumentEnvironment{readingbox}{O{}}{%
  \calloutstart{boxC}{remtitle}{Reading advice}{#1}{{\color{remtitle}\sffamily\bfseries\ ---\ }}%
}{\end{tcolorbox}}
\NewDocumentEnvironment{cautionbox}{O{}}{%
  \calloutstart{boxR}{crimson}{Caution}{#1}{{\color{crimson}\sffamily\bfseries\ ---\ }}%
}{\end{tcolorbox}}

% ---- TikZ styles -----------------------------------------------------------
\tikzset{
  blk/.style={draw=navy!70,fill=white,rounded corners=2pt,minimum height=7mm,inner sep=3pt,
              font=\sffamily\scriptsize,align=center,line width=0.6pt},
  blkfill/.style={blk,fill=mist},
  blkamber/.style={blk,draw=amber,fill=sand},
  blkteal/.style={blk,draw=teal,fill=teal!8},
  blkrose/.style={blk,draw=rose,fill=rose!8},
  flow/.style={-{Stealth[length=2.2mm,width=1.8mm]},line width=0.6pt,draw=slate},
  lab/.style={font=\sffamily\tiny,text=slate,inner sep=1pt},
}

\newcolumntype{C}{>{\centering\arraybackslash}X}

% ---- misc macros -----------------------------------------------------------
\DeclareMathOperator*{\argmin}{arg\,min}
\DeclareMathOperator*{\argmax}{arg\,max}
\newcommand{\E}{\mathbb{E}}
\newcommand{\R}{\mathbb{R}}
\newcommand{\tick}{\textcolor{teal}{\ding{51}}}
\newcommand{\cross}{\textcolor{rose}{\ding{55}}}
\newcommand{\sref}[1]{\hyperref[#1]{\S\ref*{#1}}}

% ---------------------------------------------------------------------------
%  Extras: lettered parts, units, circuit diagrams, plots, code listings and
%  four more boxes (Example, Build it, Connection, Check yourself).
%  The Key idea, Reading advice and Caution boxes are defined in preamble.tex.
% ---------------------------------------------------------------------------

% ---- lettered parts: chapters numbered C1, C2, ... ---------------------------------------
% Use a lettered part to insert a block of chapters (a crash course, say) without renumbering
% the chapters, labels and file names of the rest of the book:
%   \letterpartbegin{C} \part{...} \input{chapters/cc01_...} ... \letterpartend
% tools/build_chapter.sh numbers such files correctly when tools/book.env maps their prefix (cc:C).
\newcounter{savedchapter}
\newcommand{\letterpartbegin}[1]{%
  \setcounter{savedchapter}{\value{chapter}}\setcounter{chapter}{0}%
  \renewcommand{\thechapter}{#1\arabic{chapter}}\def\theHchapter{#1\arabic{chapter}}}
\newcommand{\letterpartend}{%
  \setcounter{chapter}{\value{savedchapter}}%
  \renewcommand{\thechapter}{\arabic{chapter}}\def\theHchapter{\arabic{chapter}}}

% ---- units (siunitx) ----------------------------------------------------------------
\sisetup{per-mode=symbol,inter-unit-product=\ensuremath{{}\cdot{}},group-separator={\,},
         group-minimum-digits=4,group-digits=integer,output-decimal-marker={.}}

% ---- circuit diagrams (circuitikz, American symbols), only if installed ------------------
% usage:  \begin{circuitikz}[ckt] ... \end{circuitikz}
% Polarity trap: to[battery1] and the source symbols put + at the START of the path; add `invert` or draw the
% path the other way, and check that current arrows leave the + terminal.
\IfFileExists{circuitikz.sty}{%
  \usepackage[american]{circuitikz}%
  \ctikzset{bipoles/length=1.15cm}}{}
\tikzset{ckt/.style={line width=0.7pt,font=\sffamily\footnotesize},
         cur/.style={-{Stealth[length=2mm,width=1.6mm]},line width=0.9pt,draw=rose},
         node dot/.style={circle,fill=black,inner sep=0.9pt}}

% ---- plots (pgfplots) -------------------------------------------------------------------
% usage:  \begin{axis}[bookplot, xlabel={...}, ylabel={...}] ... \end{axis}
% Every curve is computed from the formula that the text uses, with the numbers of the worked example.
\pgfplotsset{bookplot/.style={width=0.82\linewidth,height=5.2cm,axis lines=left,
  tick label style={font=\sffamily\scriptsize},label style={font=\sffamily\footnotesize},
  legend style={font=\sffamily\scriptsize,draw=none,fill=none},
  every axis plot/.append style={line width=0.9pt},
  title style={font=\sffamily\footnotesize}}}

% ---- code listings (ASCII only inside listings) -----------------------------------------
\lstdefinestyle{bookcode}{basicstyle=\ttfamily\footnotesize,columns=fullflexible,keepspaces=true,
  keywordstyle=\bfseries\color{maroon},commentstyle=\itshape\color{teal},stringstyle=\color{amber},
  showstringspaces=false,breaklines=true,frame=l,framerule=1.6pt,rulecolor=\color{navy!35},
  xleftmargin=1.3em,framexleftmargin=0.7em,aboveskip=6pt,belowskip=6pt,numbers=none}
\lstset{style=bookcode}

% ---- boxes: Example (rust); Build it, Connection and Check yourself (green bar) -----
%   \begin{workedbox}[title] ... \end{workedbox}        worked example with numbers
%   \begin{buildbox}[title] ... \end{buildbox}          something to build, run or measure
%   \begin{connectionbox}[title] ... \end{connectionbox} why the book's main topic cares (title in bookmeta.tex)
%   \begin{checkbox} ... \end{checkbox}                 three short questions, answers in \footnotesize
\NewDocumentEnvironment{workedbox}{O{}}{%
  \calloutstart{boxB}{exborder}{Example}{#1}{{\color{exborder}\sffamily\bfseries\ ---\ }}%
}{\end{tcolorbox}}
\NewDocumentEnvironment{buildbox}{O{}}{%
  \calloutstart{boxC}{remtitle}{Build it}{#1}{{\color{remtitle}\sffamily\bfseries\ ---\ }}%
}{\end{tcolorbox}}
\NewDocumentEnvironment{connectionbox}{O{}}{%
  \calloutstart{boxC}{remtitle}{\connectionname}{#1}{{\color{remtitle}\sffamily\bfseries\ ---\ }}%
}{\end{tcolorbox}}
\NewDocumentEnvironment{checkbox}{O{}}{%
  \calloutstart{boxC}{remtitle}{Check yourself}{#1}{{\color{remtitle}\sffamily\bfseries\ ---\ }}%
}{\end{tcolorbox}}

% ---- small math helpers ------------------------------------------------------------------
\newcommand{\transp}{^{\top}}
\newcommand{\dd}{\mathrm{d}}
     % worked/build/connection/check boxes, units, plots, code listings, lettered parts

% hyphenation help for the book's own vocabulary, e.g. \hyphenation{data-set fine-tun-ing}
